\section{Sparse matrices}
\index{matrices!sparse}

A sparse matrix is filled with mostly zeroes. Algorithms designed for sparse
matrices can avoid work by skipping operations on zeroes.

\subsection{Storage}

Sparse matrices save space and compute by only storing and operating on the
nonzero elements of the matrix. There are several possible storage
schemes\footnote{\url{https://en.wikipedia.org/wiki/Sparse_matrix\#Storage}},
but Compressed Sparse Column format is popular.

\subsection{Sparsity pattern}
\index{matrices!sparsity pattern}

The \textit{sparsity pattern} of a matrix is its pattern of structural nonzero
elements, which represents input-output coupling. An element is
\textit{structurally nonzero} if coupling exists (i.e., it \textit{can} be
nonzero), and it's \textit{structurally zero} otherwise.

\textit{Banded sparsity} is common in trajectory optimization problems.
\begin{equation*}
  \begin{bmatrix}
    \times & \times & \cdot & \cdot & \cdot & \cdot \\
    \times & \times & \cdot & \cdot & \cdot & \cdot \\
    \cdot & \cdot & \times & \times & \cdot & \cdot \\
    \cdot & \cdot & \times & \times & \cdot & \cdot \\
    \cdot & \cdot & \cdot & \cdot & \times & \times \\
    \cdot & \cdot & \cdot & \cdot & \times & \times
  \end{bmatrix}
\end{equation*}

\textit{Blocked sparsity} is common in finite element analysis problems.
\begin{equation*}
  \begin{bmatrix}
    \times & \times & \cdot & \cdot & \times & \times \\
    \times & \times & \cdot & \cdot & \times & \times \\
    \cdot & \cdot & \cdot & \cdot & \cdot & \cdot \\
    \cdot & \cdot & \cdot & \cdot & \cdot & \cdot \\
    \times & \times & \cdot & \cdot & \cdot & \cdot \\
    \times & \times & \cdot & \cdot & \cdot & \cdot
  \end{bmatrix}
\end{equation*}

When solving these two kinds of problems via iterative algorithms, the contents
of the matrix often change but the sparsity pattern is fixed, even if an element
happened to be zero for the current iteration.
\begin{remark}
  It can be beneficial for performance to keep nonstructural zeroes in the
  matrix storage if it maintains the same sparsity pattern between iterations.
\end{remark}

\subsection{Solving sparse linear systems}

Solving linear systems with a decomposition followed by a triangular solve is
more efficient and numerically stable than explicitly forming the inverse of the
system matrix. Which decomposition to use depends on the system's properties.
In general, use:
\begin{itemize}
  \item LLT for symmetric positive definite systems
  \item simplicial LDLT for symmetric positive semidefinite systems with a
    banded sparsity pattern
  \item supernodal LDLT for symmetric positive semidefinite systems with a
    blocked sparsity pattern
  \item LU for square invertible systems
  \item QR for least-squares solving
\end{itemize}

See
\url{https://libeigen.gitlab.io/eigen/docs-nightly/group__TopicSparseSystems.html}
for more specific guidance.

\textit{Pivoting} is the act of permuting rows and columns during the solve.
Factorizing a sparse matrix can result in \textit{fill-in}, where the
factorization has more nonzeros than the original matrix. For sparse matrices,
we typically want a fill-in reducing permutation.

Sparse system solvers typically have two steps: a symbolic factorization that
determines the sparsity pattern and finds a fill-in reducing permutation, and a
numerical factorization step that operates on the specific matrix. Keeping the
sparsity pattern the same between solves allows reusing the result of the first
step.
